%%%%%%%%%%%%%%%%%%%%%%%%%%%%%%%%%%%%%%%%%
% University/School Laboratory Report
% LaTeX Template
% Version 4.0 (March 21, 2022)
%
% This template originates from:
% https://www.LaTeXTemplates.com
%
% Authors:
% Vel (vel@latextemplates.com)
% Linux and Unix Users Group at Virginia Tech Wiki
%
% License:
% CC BY-NC-SA 4.0 (https://creativecommons.org/licenses/by-nc-sa/4.0/)
%
%%%%%%%%%%%%%%%%%%%%%%%%%%%%%%%%%%%%%%%%%

%----------------------------------------------------------------------------------------
%	PACKAGES AND DOCUMENT CONFIGURATIONS
%----------------------------------------------------------------------------------------

\documentclass[
	letterpaper, % Paper size, specify a4paper (A4) or letterpaper (US letter)
	11pt, % Default font size, specify 10pt, 11pt or 12pt
]{CSUniSchoolLabReport}

\usepackage{subcaption}
\usepackage{booktabs}
\geometry{top=2cm} % Override top margin to reduce space before title

%----------------------------------------------------------------------------------------
%	REPORT INFORMATION
%----------------------------------------------------------------------------------------

\title{Lab 2 - Electric Flux  \\ ECE106} % Report title

\author{Yiran \textsc{Liu} \&  Zijin \textsc{Yang}} % Author name(s), add additional authors like: '\& James \textsc{Smith}'

\date{\today} % Date of the report

%----------------------------------------------------------------------------------------

\begin{document}

\maketitle % Insert the title, author and date using the information specified above

\begin{center}
	\begin{tabular}{l r}
      Lab Section: 206 \\
      Student Number: Y.L. 21184901 \& Z.Y. 21196273
	\end{tabular}
\end{center}
%----------------------------------------------------------------------------------------
%	COMPUTATIONS
%----------------------------------------------------------------------------------------

\section{Pre-Lab: Theory}

\subsection{Eletric Flux through Closed Sphereical Surface}

A point charge of $+0.75\mu C$ is placed at $(x, y, z) = (0, 0, 0)$, we are to compute the total eletric flux through a sphereical surface of radius $R=2.0m$ from the definition of eletric flux. 

\[
  \Phi_E = \iint \vec{E} \cdot d\vec{S} = \int_{0}^{2 \pi} \int_{0}^{2 \pi} \frac{1}{4 \pi \epsilon_0} \frac{q}{R^2} \hat{r} \cdot \hat{E} \ r d\theta \ r d\phi
\]
where $\epsilon_{0} = 8.854187 \times 10^{-12} \ F / m$ and $q = +0.75 \times 10^{-6} C$. 

We notice that the eletric field strength $\vec{E}$ is the same, and that $\hat{r} \cdot \hat{E} = 1$. Therefore, we can compute the flux as:
\[
  \Phi_E = S \|\vec{E}\| = 4 \pi R^2 \frac{q}{4 \pi R^2 \epsilon_0} 
\]
We compute for $R=2.0m$ and $R=5.0m$, $\Phi_{2m} = \Phi_{5m} = 84705.69 Nm^2/C$.

\subsection{Charge Density}

The charge density of a spherical charge cloud with radius $r=0.01m$ and charge of $+0.75\mu C$. The volume of the sphere is $V = \frac{4}{3} \pi r^3 = 4.1888 \times 10^{-6}$, density $\rho=\frac{q}{V} = 0.17905 C/m^3$

\subsection{Eletric Flux Through Square Surfaces}

We denote the the eletric flux through the three squares as $\Phi_1$, $\Phi_2$ and $\Phi_3$, respectively. We construct cubic gaussian surfaces for square 1 and 2, denote as Cube1 and Cube2. Using the Gauss' Law and considering the symetry:

\[
4\Phi_1 = 4\Phi_2 = \frac{q_\text{enclosed}}{\epsilon_0}
\]

We observe that Square 3 contains all of Square 1 and some additional area that has non-zero eletric flux. Therefore, $\Phi_3 > \Phi_1$. The final relationship is:

\[
\Phi_3 > \Phi_2 = \Phi_1
\]

\section{Simulation for Spherical Gaussian Surfaces}


\end{document}